\chapter*{Conclusions}
\phantomsection
\addcontentsline{toc}{chapter}{Conclusions}

This project set out to design and implement a domain-specific language (DSL) to partially automate and standardize cybersecurity incident reporting. To achieve the goals outlined in the Abstract, the following objectives were completed:

\begin{itemize}
    \item \textbf{Formal Specification:} A custom grammar (\texttt{IncidentFlow.g4}) was implemented using ANTLR, defining language constructs for triggers, phases, logic branching, and severity escalation.
    \item \textbf{Compiler and Validation:} A parser pipeline was established to evaluate the incident playbooks for both syntactic correctness and semantic consistency, effectively replacing manual requirement-checking with automated validation.
    \item \textbf{CLI Toolkit:} A command-line interface was developed to allow engineers to validate scripts and generate structured Markdown incident reports. Specifically, an interactive scenario wizard was added to accelerate playbook generation for common MITRE ATT\&CK classifications.
    \item \textbf{API and Frontend Integration:} The core compiler was exposed via an HTTP REST API, which served as the backend for a React and Vite-based web frontend. This enabled real-time syntax checking and report rendering directly in the browser.
\end{itemize}

The resulting implementation confirms that replacing manual, document-centric incident reporting with a process-centric DSL is technically feasible. The unified architecture successfully bridges the operational necessity of executing technical responses with the administrative requirement of maintaining auditable compliance evidence. 

Future development could extend the DSL's capabilities by directly integrating it with enterprise SOAR systems, allowing the compiler to not only generate documentation but simultaneously push automated mitigation commands across network endpoints.

The IncidentFlow DSL project source code relative to this implementation is available at: \url{https://github.com/nnorian/IncidentFlow_dsl/} \autocite{report2026}.
